\chapter{ตารางจำแนกจุลินทรีย์ในพืชเศรษฐกิจไทย}
\label{app:economic_crops}

\begin{table}[H]
\centering
\small
\begin{tabularx}{\textwidth}{|p{2.6cm}|X|X|}
\hline
\rowcolor{agriMint}
\textbf{พืชเศรษฐกิจ} & \textbf{จุลินทรีย์ที่เป็นประโยชน์} & \textbf{จุลินทรีย์ก่อโรคสำคัญ} \\ \hline
\textbf{ข้าว (Rice)} & 
- \textit{Azospirillum brasilense} (ตรึงไนโตรเจน)\newline
- \textit{Rhodopseudomonas palustris} (ลดซัลไฟด์)\newline
- \textit{Trichoderma asperellum} & 
- \textit{Magnaporthe oryzae} (โรคไหม้)\newline
- \textit{Rhizoctonia solani} (โรคกาบใบแห้ง)\newline
- \textit{Xanthomonas oryzae} (โรคขอบใบแห้ง) \\ \hline

\textbf{ทุเรียน (Durian)} & 
- \textit{Trichoderma harzianum} (ชีววิธี)\newline
- \textit{Bacillus amyloliquefaciens}\newline
- \textit{Rhizophagus irregularis} (ไมคอร์ไรซา) & 
- \textit{Phytophthora palmivora} (รากเน่าโคนเน่า)\newline
- \textit{Colletotrichum gloeosporioides} (แอนแทรคโนส)\newline
- \textit{Fusarium solani} (เน่าร่วม) \\ \hline

\textbf{ยางพารา (Rubber)} & 
- เชื้อราไมคอร์ไรซา (AMF)\newline
- \textit{Streptomyces} spp. (ยับยั้งเชื้อรา) & 
- \textit{Phytophthora botryosa} (ใบร่วง)\newline
- \textit{Rigidoporus microporus} (โรครากขาว)\newline
- \textit{Corynespora cassiicola} (ใบร่วง) \\ \hline

\textbf{มันสำปะหลัง} & 
- \textit{Glomus} spp. (ช่วยดูดฟอสฟอรัส)\newline
- \textit{Bacillus subtilis} (กระตุ้นแตกราก) & 
- \textit{Xanthomonas axonopodis} (โรคใบไหม้)\newline
- ไวรัสใบด่างมันสำปะหลัง (SLCMV) \\ \hline

\textbf{อ้อย (Sugarcane)} & 
- \textit{Gluconacetobacter diazotrophicus}\newline
  (แบคทีเรียเอนโดไฟต์ตรึงไนโตรเจนสูง) & 
- \textit{Sporisorium scitamineum} (โรคแส้ดำ)\newline
- ไฟโตพลาสมา (โรคใบขาวอ้อย) \\ \hline
\end{tabularx}
\caption{การกระจายตัวของจุลินทรีย์ที่มีประโยชน์และจุลินทรีย์ก่อโรคในพืชเศรษฐกิจหลักของไทย}
\label{tab:crop_microbes}
\end{table}

\clearpage